% Complete primary and fresh decoded native genus67 RR checks.
\subsection{Excluding a point of order $69$}
\begin{proposition}\label{prop:decic-no69}
An elliptic curve over a number field of degree $10$ has no point of order $69$.
\end{proposition}
\begin{proof}
Put $X=X_1(69)$, of genus $133$. By \cite[Theorem~3.1]{DEHMZ},
$\operatorname{rank}J_1(69)(\mathbb Q)=0$.
The gonality bound gives
\[
 \operatorname{gon}_{\overline{\mathbb Q}}X\geq
 \frac{325\cdot2112}{32768}=\frac{10725}{512}>20.
\]
Suppose that a proposed torsion point gives a rational effective
noncuspidal divisor $D$ of degree $10$ on $X$.

The cusp signatures over $\mathbb Q$ and $\mathbb F_2$ are
\[
 1^{22},2^{11},22^2\qquad\text{and}\qquad1^{22},2^{11},11^2,22.
\]
Every residue cusp of degree at most $10$ is therefore an intact
global Galois orbit. If $\overline D$ is entirely cuspidal, its
cuspidal lift $C$ is a rational effective divisor of degree $10$.
The torsion kernel of reduction at $2$ has exponent two, so
$2[D-C]=0$. Indeed, multiplication by two sends the first formal
subgroup into the torsion-free second formal subgroup.
Consequently $2D\sim2C$. The gonality bound forces $D=C$, a
contradiction.

The noncuspidal closed-point counts on $X_{\mathbb F_2}$ in
degrees $1$ through $10$ are
\[
 (0,0,0,0,0,0,66,44,44,44).
\]
They are independently verified from every finite elliptic-curve
isomorphism class and its level structures. Hence a reduction
with noncuspidal support contains one point of degree
$e\in\{7,8,9,10\}$ and a cuspidal complement of degree $10-e$.

First use the quotient
\[
 Y=X/\ker\varepsilon,\qquad
 \varepsilon(a)=\left(\frac{69}{a}\right).
\]
It is an unramified double cover of $X_0(69)$, of genus $13$.
Both quotient Jacobians have rank zero. Since
$\varepsilon(5)=1$, the operator $T_5-6$ kills their rational
Jacobians. The base model is the normalization of
$X_0(23)\times_{X(1)}X_0(3)$: over the genus-two field of
$X_0(23)_{\mathbb F_2}$, with its exact $j$-map, adjoining a root of
\[
 v^4+v^3+1/j=0
\]
gives the genus-seven field of $X_0(69)$.
The double cover is identified by its unramified quadratic
character. The base Picard group is
$C_{22}\times C_{44}\times\mathbb Z$.
Its prescribed cusp Frobenius values leave two characters;
requiring the three actual degree-seven source images to split
leaves one. Exact field, cusp and Picard-generator checks identify
this cover with $Y_{\mathbb F_2}$.

The complete base and lifted Hecke tests give the following table.
The lift column counts all choices above the base Hecke divisor,
so it includes the actual correspondence.
\[
\begin{array}{c|r|r|r|r}
 e&\text{base points}&\text{nonprincipal base tests}
 &\text{lift choices tested}&\text{principal lift choices}\\\hline
 7&3&2&8&0\\
 8&2&2&0&0\\
 9&2&1&2&0\\
 10&2&0&8&2
\end{array}
\]
All $27$ degree-zero divisor tests are replayed in fresh native
maximal orders. Every serialized prime uses two generators,
with equality of the generated ideal and exact function
roundtrips. Independent explicit computations reproduce the
$54$ underlying $5$-isogenies, the $23$-isogeny coordinates, and
all norm multiplicities.
The cuspidal complements contribute zero: they are sums of
complete global cusp orbits of degrees one and two, and the
cusp Hecke formula together with $\varepsilon(5)=1$ gives
$T_5\pi_*C'=6\pi_*C'$.
Thus only one degree-ten source diamond orbit remains.
The two principal choices in the last row lie above that same
orbit; they are not asserted to be the actual Hecke divisor.

We resolve it by retaining the generator of the $23$-part.
Let
\[
 Z=\operatorname{Norm}\bigl(X_1(23)\times_{X(1)}X_0(3)\bigr).
\]
This is the degree-two quotient $X/\langle\langle22\rangle\rangle$,
of genus $67$, parametrizing $(E,\pm P_{23},C_3)$.
Its Jacobian has rank zero, being a quotient of $J_1(69)$.
Over $K=\mathbb F_2(X_1(23))$, of genus $12$, its field is
\[
 K(v),\qquad v^4+v^3+1/j=0.
\]
The quartic is irreducible and has genus $67$. The family
$y^2+xy+vy=x^3$, with its point $(0,0)$ of order three, verifies
the formula $j=1/(v^3(v+1))$ and hence the modular identification.
At a rational base cusp, $v=1$ is a simple root of the reduced
quartic: $f(1)=1/j$ has positive valuation and $f'(1)=1$.
Hensel's lemma gives a rational place, proving that the exact
constant field is $\mathbb F_2$.

The remaining orbit has
\[
 j^{10}+j^8+j^7+j^6+j^5+j^4+j^3+j+1=0.
\]
The qualifying curve over $\mathbb F_{1024}$ has cyclic group of
order $966$. Its subgroup of order $69$ is unique, so all its
level structures form one diamond orbit. Fix a generator $P_{69}$.
We construct its image on $Z$ from $(E,3P_{69})$ on $X_1(23)$.
At this degree-ten base point, the quartic has exactly two
irreducible factors, of degrees one and three. The unique linear
factor therefore identifies the actual source image on $Z$.

The six $5$-isogenies form two Frobenius orbits, with one and five
members. Their images on $X_1(23)$ have degrees $10$ and $50$,
equal to their $j$-degrees and their full $69$-level degrees.
At each image the quartic again factors as $1+3$.
Thus each actual Hecke image has a unique lift of the same
degree to $Z$. This removes the ambiguity present in the
smaller double cover. The norm weights are
$10\cdot1/10=1$ and $10\cdot5/50=1$, so the resulting effective
Hecke divisor $T$ has degree $60$.
The image of $5P_{69}$ is constructed in the same way.
All prime restrictions are checked with both generators of the
base prime and a uniformizer. The proof packet records two
actual generators of every prime, verifies their generated
ideals and all serialization roundtrips, and reconstructs the
source, diamond and Hecke divisors exactly.

Write $P_Z$ for the degree-ten source image. Exact Riemann--Roch,
independently repeated on the decoded packet in a fresh model,
gives
\[
 h^0\bigl(T-P_Z-5\langle5\rangle P_Z\bigr)=0.
\]
In particular the degree-zero divisor is nonprincipal.
On the other hand, Eichler--Shimura at $5$ and injectivity of
rational torsion reduction at the good odd prime $5$ show that
$T_5-1-5\langle5\rangle$ kills $J_Z(\mathbb Q)$.
At the image $c$ of the standard rational Tate cusp,
$T_5c=c+5\langle5\rangle c$ as a literal divisor.
The image of $D$ on $Z$ would therefore give a principal divisor
with exactly the displayed reduction, a contradiction.
Hecke operators commute with diamonds, so this excludes the
entire remaining source orbit and completes the proof.
\end{proof}
